\chapter{Literature Review}
\section{General}
The seismic performance of reinforced concrete buildings has been widely studied because of the vulnerability of reinforced concrete (RC) frame structures during major earthquakes. In seismic regions, the behaviour of a building depends not only on member strength but also on stiffness distribution, ductility, structural regularity, detailing quality, material properties, and construction practice. For low-rise reinforced concrete buildings in Nepal, simplified design and detailing provisions such as the Mandatory Rules of Thumb (MRT) and Ready-to-Use Detailing (RUD) guidelines are commonly used. The Draft Final NBC 205:2012 provides an intermediate ready-to-use reference between the earlier MRT document and the current NBC 205:2024 guideline \cite{nbc205_2012}. These guidelines are intended to improve minimum seismic safety in ordinary residential buildings, especially in locations where detailed engineering design and professional supervision may not always be available \cite{nbc205_1994,nbc205_2024}.
Past earthquake events in Nepal have shown that many residential buildings are vulnerable due to poor construction quality, insufficient ductile detailing, weak-column strong-beam behaviour, irregular configuration, and inadequate implementation of building code provisions \cite{chaulagain2013seismic,parajuli2015gorkha}. Therefore, it is necessary to review previous studies related to Nepalese RC buildings, MRT/RUD-based design, structural irregularity, torsional behaviour, and performance-based seismic evaluation. This literature review presents previous research related to these themes and establishes the basis for the methodology adopted in the present study.
\section{Seismic Vulnerability of RC Buildings in Nepal}
Reinforced concrete frame buildings have become one of the most common forms of construction in Nepal, particularly in urban and semi-urban areas. However, many of these buildings are constructed by owner-builders with the involvement of local technicians and masons rather than through complete engineering design and supervision. As a result, several buildings may lack adequate seismic detailing and may not satisfy modern performance requirements.
Chaulagain et al. \cite{chaulagain2013seismic} studied the seismic response of current reinforced concrete buildings in Nepal by comparing different design and construction practices. Their study considered different categories of building design and construction practice and compared their seismic response using nonlinear static and dynamic analyses. The results were compared in terms of base shear, roof displacement, capacity curve, and inter-storey drift. The study showed that buildings representing current construction practice and older code provisions experienced higher drift demands and greater vulnerability than well-designed buildings. This work is important for the present research because it shows the need to evaluate Nepalese RC buildings using nonlinear performance-based methods rather than relying only on prescriptive design requirements.
The seismic vulnerability of Nepalese RC buildings is also related to construction practice and code implementation. Many low-rise residential buildings are constructed with RC frames and masonry infill walls, with inadequate consideration of seismic detailing. Although building codes provide minimum requirements, improper workmanship, lack of supervision, poor material quality, and irregular architectural layouts can significantly affect actual seismic performance. Therefore, the evaluation of typical Nepalese low-rise RC buildings remains an important research area.

\begin{figure}[!htbp]
\centering
\includegraphics[width=0.96\textwidth]{figures/ch2/seismic_vulnerability_rc_building.png}
\caption{Seismic vulnerability of low-rise RC buildings.}
\label{fig:rc_vulnerability_concept}
\end{figure}

Figure~\ref{fig:rc_vulnerability_concept} summarizes the main vulnerability mechanisms discussed in the literature for low-rise reinforced-concrete buildings. Inadequate detailing, workmanship and material quality can increase damage susceptibility, while plan irregularity and eccentricity can disturb stiffness and force distribution and increase torsional and drift concentration. The schematic uses the three-storey study context and an illustrative re-entrant plan based on the study grid only to communicate these mechanisms; it is not an analysis result or a fragility curve. This interpretation is consistent with the Nepalese RC-building observations and irregular-building studies reviewed in this chapter \cite{chaulagain2013seismic,parajuli2015gorkha,poudel2017seismic}.
\FloatBarrier

\section{MRT and RUD-Based Design Provisions}
The Mandatory Rules of Thumb were introduced in Nepal to provide simplified design and detailing requirements for ordinary low-rise reinforced concrete residential buildings. These rules provide guidance regarding member dimensions, reinforcement details, and general restrictions for ordinary buildings \cite{nbc205_1994}. The Draft Final NBC 205:2012 document represents a ready-to-use detailing stage for low-rise RC buildings without masonry infill and provides the basis for the representative existing-building evaluation in Phase I \cite{nbc205_2012}. The current NBC 205:2024 Ready-to-Use Detailing guideline provides updated provisions for the regular verification model in Phase II and the irregular extensions in Phase III, while NBC 105:2025 provides the seismic demand definitions \cite{nbc205_2024,nbc105_2025}.

Sthapita and Shrestha~[26] provide Nepal-specific context for the reassessment of earlier ready-to-use provisions under updated seismic demand. They analysed three low-rise RC building samples based on NBC 205:1994, NBC 201:1994, and NBC 205:2012 using ETABS under NBC 105:2020. Their results indicated that the 1994 MRT samples showed several deficiencies under the updated seismic loading, while the 2012 sample performed comparatively better in storey drift and deflection although reinforcement and ductile-detailing deficiencies remained. This supports rechecking legacy ready-to-use buildings when the seismic-demand framework changes and provides relevant context for the present comparison of the Draft Final NBC 205:2012 reference with the NBC 205:2024 benchmark.
Poudel et al. \cite{poudel2025seismic} evaluated the seismic performance of low-rise reinforced concrete framed buildings designed using the Ready-to-Use Detailing guideline NBC 205:2024 in Nepal. Their study used nonlinear pushover analysis to assess capacity curves, storey drift, base shear, performance points, and hinge formation. The buildings were analysed under different seismic demands and soil types. The study concluded that buildings designed according to the updated RUD guidelines may provide adequate seismic performance when recommended material quality, detailing, and construction standards are satisfied. This study is directly relevant to the present research because it evaluates NBC 205:2024-based low-rise RC buildings using performance-based analysis.
However, most simplified guidelines are primarily suitable for regular and ordinary building configurations. In practice, many residential buildings deviate from regular rectangular plans due to land shape, functional requirements, architectural preference, and space limitations. Therefore, the adequacy of MRT/RUD-based provisions for plan-irregular buildings requires further study. The present research extends this idea by evaluating NBC 205:2024 ready-to-use RC buildings with different plan-irregular configurations.
\section{MRT Non-Compliance and Geometric Limitations}
The seismic performance of MRT-based buildings is affected by whether the buildings remain within the geometric limitations prescribed by the guidelines. Poudel and Shankar \cite{poudel2017seismic} studied the seismic vulnerability of RC frame buildings not complying with MRT provisions based on slenderness ratio and plan aspect ratio. Their study considered buildings exceeding MRT limitations and evaluated seismic vulnerability using pushover analysis, linear time-history analysis, and fragility curves. The results showed that buildings exceeding the recommended limits may experience higher vulnerability and increased probability of damage.
This study is important because it indicates that even if a building follows simplified design concepts, geometric irregularity or non-compliance with prescribed limitations can significantly affect seismic performance. The present study is related to this issue because plan and height irregularity can disturb the expected behaviour of RUD-designed buildings. While Poudel and Shankar focused on slenderness ratio and plan aspect ratio, the present Phase III assessment evaluates L-type, L1-type, L2-type, T-type, U-type, and Non-Parallel three-storey buildings and treats a separate four-storey-plus-stair-cover model as a height-extension case.
\section{Structural Irregularity and Seismic Response}
Structural irregularity is one of the major causes of undesirable seismic behaviour in buildings. Irregularity may occur due to discontinuity or non-uniform distribution of mass, stiffness, strength, or geometry. A regular structure generally has uniform distribution of stiffness and mass in both plan and elevation, whereas irregular structures may show abrupt changes or asymmetry. During earthquakes, irregularity can cause concentration of forces and deformation in certain members or storeys.
Plan irregularity is particularly important because it can produce torsional response. When the centre of mass and centre of rigidity do not coincide, lateral seismic forces may generate torsional moments. This causes some parts of the building to experience larger displacement and drift than others. Re-entrant corners, L-shaped plans, T-shaped plans, U-shaped plans, and non-parallel framing arrangements are common examples of plan irregularity. Such configurations can alter the lateral load path and create uneven distribution of seismic demand.
Jadhav and Shelar \cite{jadhav2025study} discussed the dynamic analysis of irregular shape structures using time-history and response spectrum analysis. Their study emphasized that irregularity in mass, stiffness, and geometry affects the dynamic characteristics of buildings. The study also highlighted the importance of using structural analysis software such as ETABS for evaluating displacement and other seismic response parameters in irregular buildings. This supports the use of response spectrum analysis and ETABS modelling in the present study.
\section{Torsional Behaviour of Irregular Buildings}
Torsional response is a key concern in plan-irregular buildings. In a regular structure, the centre of mass and centre of rigidity are generally close to each other, resulting in predominantly translational response under lateral loading. However, in irregular buildings, the eccentricity between the centre of mass and centre of rigidity may become significant. This eccentricity causes CM--CR eccentricity of the floor diaphragm and increases seismic demand at the edges of the building.
The evaluation of torsional irregularity is therefore necessary for understanding the seismic behaviour of plan-irregular structures. Important parameters include centre of mass coordinates, centre of rigidity coordinates, eccentricity in the principal directions, storey displacement, and inter-storey drift. Rigid diaphragm modelling is commonly used to ensure proper transfer of lateral forces and to evaluate torsional response at each floor level.
Cimellaro et al. \cite{cimellaro2014bidirectional} studied bidirectional pushover analysis of irregular structures. Their research discussed the limitations of conventional pushover methods when applied to irregular buildings subjected to bidirectional ground motions. They found that torsional effects, higher mode contribution, and bidirectional seismic action can significantly influence the response of irregular buildings. Their work is useful for the present study because it supports the need to carefully interpret pushover results for plan-irregular and torsionally sensitive buildings.
\section{Equivalent Static Method}
The Equivalent Static Method is one of the simplest methods for seismic analysis of buildings. In this method, earthquake effects are represented by equivalent lateral forces applied along the height of the structure. The method involves determining seismic weight, calculating design base shear, distributing lateral forces vertically, and analysing the resulting displacement, drift, and storey shear. It is commonly used for regular and low-rise buildings where higher mode effects are not dominant.
In the context of the present study, equivalent static analysis is useful for obtaining initial estimates of lateral force demand and for comparing the response of regular and irregular buildings under different seismic zone factors and soil conditions. It also provides a reference for comparing dynamic response spectrum results. However, for irregular buildings, equivalent static analysis alone may not fully capture torsional effects and modal behaviour. Therefore, it is supplemented by response spectrum analysis and nonlinear pushover analysis.
\section{Response Spectrum Analysis}
Response Spectrum Analysis is a linear dynamic method widely used to evaluate the seismic response of buildings. It considers the contribution of different vibration modes and combines modal responses to estimate structural demand. This method is more suitable than the equivalent static method for buildings where dynamic characteristics, modal participation, or irregularity influence the response.
In seismic design, response spectrum analysis is commonly used to obtain natural time periods, frequencies, mode shapes, modal mass participation ratios, base shear, storey displacement, and inter-storey drift. In the present study, response spectrum functions based on NBC provisions are used for different soil types and seismic zone factors. The results obtained from response spectrum analysis help evaluate the influence of plan irregularity on dynamic characteristics and elastic seismic response.
NBC 105:2025 provides the seismic design basis for Nepal, including response spectrum parameters, seismic zone factors, importance factor, and soil classification \cite{nbc105_2025}. Since the present study considers Soil Type C, Soil Type D, and different seismic zone factors, response spectrum analysis is an essential part of the methodology.
\section{Nonlinear Static Pushover Analysis}
Nonlinear static pushover analysis is widely used for performance-based seismic evaluation of buildings. In pushover analysis, the structure is subjected to gradually increasing lateral loads until a target displacement or failure condition is reached. The resulting relationship between base shear and roof displacement is known as the capacity curve. This curve provides insight into the strength, stiffness, ductility, and deformation capacity of the structure.
ATC-40 \cite{atc1996atc40}, FEMA 356 \cite{fema2000fema356}, FEMA 440 \cite{fema2005fema440}, and ASCE 41-17 \cite{asce2017asce41} provide guidelines for nonlinear static evaluation and performance-based assessment. These documents define procedures for estimating performance points, evaluating deformation demand, and classifying structural performance levels. Performance levels such as Immediate Occupancy, Life Safety, and Collapse Prevention are commonly used to interpret the expected damage state of a structure.
Poudel et al. \cite{poudel2025seismic} used nonlinear pushover analysis to evaluate low-rise RC buildings designed according to NBC 205:2024. Their work demonstrated the usefulness of pushover analysis for studying capacity curves, performance points, storey drift, and hinge formation. Similarly, Chaulagain et al. \cite{chaulagain2013seismic} used nonlinear static and dynamic analyses to evaluate the seismic response of different RC building design practices in Nepal. These studies support the use of pushover analysis in the present research.
\section{Performance Point and Capacity Curve}
The performance point represents the estimated seismic demand on the structure and is commonly obtained by comparing the structural capacity spectrum with the seismic demand spectrum. It gives the roof displacement and base shear at which the structure is expected to perform under a given earthquake demand. The performance point is useful for comparing different building configurations because it reflects both demand and capacity.
The capacity curve, plotted as base shear versus roof displacement, shows the progressive nonlinear response of the structure. Initially, the curve is approximately linear, representing elastic behaviour. After yielding begins, stiffness reduces and plastic hinges form in beams or columns. The shape of the capacity curve indicates the strength reserve, deformation capacity, and ductility of the structure. In performance-based evaluation, comparison of capacity curves among different building models helps identify which configuration has higher strength, stiffness, or displacement capacity.
In the present study, capacity curves and performance points are used to compare regular and plan-irregular NBC 205:2024 RC buildings. The comparison helps evaluate whether plan irregularity changes seismic capacity and performance-point response.
\section{Plastic Hinge Formation and Performance Levels}
Plastic hinge formation is an important indicator of nonlinear structural behaviour. During pushover analysis, plastic hinges form when structural members exceed their elastic capacity. The location and sequence of hinge formation help identify the probable damage pattern and failure mechanism. A desirable seismic mechanism is generally associated with hinge formation in beams before columns, which supports the strong-column weak-beam concept.
ASCE 41-17 \cite{asce2017asce41} provides acceptance criteria for nonlinear modelling and performance evaluation of structural members. Beam hinges are commonly modelled as moment hinges, while column hinges are commonly modelled using axial force and biaxial moment interaction. In the present study, the available final-step hinge states are interpreted using the Immediate Occupancy, Life Safety, and Collapse Prevention labels. Detailed hinge-assignment definitions are not restated because the complete model-definition record is not available for every case.
Paulay and Priestley \cite{paulay1992seismic} provided important concepts related to seismic design of reinforced concrete structures, ductility, and plastic hinge behaviour. Mander et al. \cite{mander1988confined} developed a stress-strain model for confined concrete, which is widely used in nonlinear modelling of RC members. Menegotto and Pinto \cite{menegotto1973cyclic} proposed a model for cyclic behaviour of steel reinforcement. These references form the theoretical basis for nonlinear material behaviour and hinge modelling.
\section{Inter-Storey Drift and Damage Assessment}
Inter-storey drift is one of the most important parameters in seismic performance evaluation. It measures the relative lateral displacement between two consecutive storeys and is closely related to structural and non-structural damage. Excessive drift can cause cracking, member yielding, damage to infill walls, and possible instability. Therefore, drift is commonly used as a performance indicator in seismic design guidelines.
ATC-40 and FEMA 356 provide drift-based performance limits for different damage states \cite{atc1996atc40,fema2000fema356}. Chaulagain et al. \cite{chaulagain2013seismic} compared inter-storey drift demands of different Nepalese RC building types and used drift limits to assess seismic safety. Their study showed that poorly detailed and current construction practice buildings may experience higher drift demands. This supports the use of drift as a key response parameter in the present research.
In this study, inter-storey drift is evaluated for all regular and plan-irregular models under different seismic zone factors and soil conditions. Drift comparison helps identify which building configuration is more flexible and which configuration may show relatively higher deformation demand under lateral seismic loading.
\section{Base Shear and Storey Shear}
Base shear represents the total lateral seismic force resisted at the base of a structure. It is one of the primary indicators of seismic demand and structural strength. In equivalent static and response spectrum analysis, base shear is used to compare lateral force demand among different models. In pushover analysis, base shear capacity is used to evaluate the maximum lateral resistance of the structure.
Changes in building configuration can affect base shear due to changes in mass, stiffness, natural period, and dynamic characteristics. Plan-irregular buildings may show different base shear response in the two principal directions due to asymmetry and different stiffness distribution. Therefore, base shear comparison is important for understanding the influence of plan irregularity on lateral force resistance.
\section{Natural Time Period and Frequency}
The natural time period and frequency of a structure describe its dynamic characteristics. The fundamental time period depends on mass and stiffness distribution. A more flexible structure generally has a longer time period and lower frequency, while a stiffer structure has a shorter time period and higher frequency. Structural irregularity can change the dynamic characteristics of a building by altering stiffness and mass distribution.
In response spectrum analysis, natural time period and modal frequency are important because seismic demand depends on the period range of the structure. If the natural frequency of the building is close to the dominant frequency of earthquake excitation, amplification of response may occur. Poudel and Shankar \cite{poudel2017seismic} discussed the importance of dynamic characteristics in explaining seismic vulnerability trends in MRT non-compliant buildings. These parameters remain relevant literature background; however, they are not used as quantitative comparison results in this study.
\section{Influence of Soil Type and Seismic Zone Factor}
Seismic demand depends on both seismic hazard level and local soil condition. NBC 105:2025 classifies soil types and provides response spectrum parameters for seismic design \cite{nbc105_2025}. Softer soil conditions may amplify seismic response and increase displacement and drift demand. Similarly, higher seismic zone factors increase design seismic demand.
Poudel et al. \cite{poudel2025seismic} evaluated RUD-NBC 205:2024 buildings for different seismic demands and soil types. Their study considered seismic zone factors ranging from 0.25 to 0.40 and different soil conditions. The present study adopts a similar range of seismic zone factors and soil conditions to evaluate how NBC 205:2024 regular and plan-irregular RC buildings respond under varying seismic demand.
\section{Summary of Major Literature Reviewed}
The major literature reviewed in this study is summarized in \mytab{tab:literature_research_papers}, \mytab{tab:literature_codes_guidelines}, and \mytab{tab:literature_nonlinear_modelling}. The reviewed literature includes previous research on Nepalese reinforced concrete buildings, MRT/RUD-based design provisions, seismic irregularity, nonlinear pushover analysis, performance-based assessment, and nonlinear modelling of reinforced concrete members.
\begingroup
\small
\setlength{\LTpre}{6pt}
\setlength{\LTpost}{6pt}
\begin{longtable}{@{}p{3.35cm}p{2.95cm}p{3.35cm}p{3.25cm}@{}}
\caption{Summary of major research papers reviewed}\label{tab:literature_research_papers}\\
\hline\hline
\textbf{Author(s)} & \textbf{Study Focus} & \textbf{Methodology Adopted} & \textbf{Relevance to Present Study} \\ \hline
\endfirsthead
\hline\hline
\textbf{Author(s)} & \textbf{Study Focus} & \textbf{Methodology Adopted} & \textbf{Relevance to Present Study} \\ \hline
\endhead
Chaulagain et al.\newline \cite{chaulagain2013seismic}
& Seismic response of current reinforced concrete buildings in Nepal under different design and construction practices.
& Nonlinear static analysis, nonlinear dynamic analysis, capacity curve, base shear, roof displacement, and inter-storey drift comparison.
& Supports the need for performance-based evaluation of Nepalese RC buildings and comparison of drift, base shear, and capacity response. \\ \hline
Poudel and Shankar\newline \cite{poudel2017seismic}
& Seismic vulnerability of RC frame buildings not complying with MRT limitations based on slenderness ratio and plan aspect ratio.
& ETABS modelling, pushover analysis, linear time-history analysis, and fragility curve development.
& Shows that MRT geometric limitations affect seismic vulnerability and supports evaluation of MRT-designed buildings with irregular configurations. \\ \hline
Sthapita and Shrestha~[26]
& Analysis of MRT-compliant buildings under NBC 105:2020.
& ETABS assessment of samples based on NBC 205:1994, NBC 201:1994, and NBC 205:2012 using drift, displacement, and member design checks.
& Provides Nepal-specific evidence that earlier MRT/RUD provisions may require reassessment under updated seismic demand and supports the comparison with the 2012 reference building. \\ \hline
Poudel et al.\newline \cite{poudel2025seismic}
& Seismic performance evaluation of low-rise RC framed buildings designed using RUD-NBC 205:2024.
& ETABS modelling, nonlinear pushover analysis, response spectrum demand, performance point, storey drift, and hinge formation assessment.
& Directly supports the present methodology involving NBC-based demand, pushover analysis, capacity curves, performance points, and hinge-state evaluation. \\ \hline
Cimellaro et al.\newline \cite{cimellaro2014bidirectional}
& Bidirectional pushover analysis of irregular structures.
& Bidirectional pushover analysis, nonlinear response history analysis, inter-storey drift, and floor rotation comparison.
& Supports discussion of torsional effects, limitations of conventional pushover analysis, and seismic response of irregular buildings. \\ \hline
Jadhav and Shelar\newline \cite{jadhav2025study}
& Dynamic analysis of irregular shape structures using time-history and response spectrum analysis.
& Study on irregular RC buildings, response spectrum analysis, time-history analysis, and ETABS modelling.
& Supports the use of dynamic analysis for irregular structures and highlights the importance of displacement and seismic response evaluation. \\ \hline\hline
\end{longtable}
\endgroup
\begingroup
\footnotesize
\renewcommand{\arraystretch}{1.08}
\setlength{\tabcolsep}{3pt}
\setlength{\LTpre}{6pt}
\setlength{\LTpost}{6pt}
\begin{longtable}{@{}p{3.25cm}p{3.30cm}p{3.35cm}p{2.90cm}@{}}
\caption{Summary of codes, standards, and seismic evaluation guidelines reviewed}\label{tab:literature_codes_guidelines}\\
\hline\hline
\textbf{Reference} & \textbf{Main Content} & \textbf{Application in Seismic Evaluation} & \textbf{Relevance to Present Study} \\ \hline
\endfirsthead
\hline\hline
\textbf{Reference} & \textbf{Main Content} & \textbf{Application in Seismic Evaluation} & \textbf{Relevance to Present Study} \\ \hline
\endhead
NBC 205:1994 \cite{nbc205_1994}
& Mandatory Rules of Thumb for low-rise reinforced concrete buildings.
& Provides simplified member dimensions and detailing provisions for ordinary RC buildings.
& Forms the historical background of simplified NBC 205 provisions used in Nepal. \\ \hline
Draft Final NBC 205:2012 \cite{nbc205_2012}
& Ready-to-Use Guideline for detailing low-rise reinforced concrete buildings without masonry infill.
& Provides the documented basis for the representative existing-building model evaluated in Phase I.
& Supports evaluation of the 2012 building stock and comparison with the revised regular 2024 benchmark. \\ \hline
NBC 205:2024 \cite{nbc205_2024}
& Ready-to-Use Detailing guideline for low-rise reinforced concrete buildings without masonry infill.
& Provides updated detailing provisions for low-rise RC buildings considering seismic zone factors and soil types.
& Supports the design basis and comparison of MRT/RUD-\allowbreak{}designed RC buildings. \\ \hline
NBC 105:2025 \cite{nbc105_2025}
& Nepal seismic design code.
& Provides seismic zone factor, soil classification, response spectrum, and seismic design requirements.
& Used as the basis for seismic parameters, response spectrum functions, and soil conditions. \\ \hline
ATC-40 \cite{atc1996atc40}
& Seismic evaluation and retrofit of concrete buildings.
& Provides capacity spectrum method, performance point evaluation, and performance-based seismic assessment.
& Supports interpretation of capacity curves and performance-based seismic evaluation. \\ \hline
FEMA 356 \cite{fema2000fema356}
& Prestandard and commentary for seismic rehabilitation of buildings.
& Provides performance levels, drift limits, and nonlinear evaluation procedures.
& Supports drift-based seismic assessment and performance-level interpretation. \\ \hline
FEMA 440 \cite{fema2005fema440}
& Improvement of nonlinear static seismic analysis procedures.
& Provides improved procedures for nonlinear static analysis and performance point estimation.
& Supports pushover-based performance point evaluation. \\ \hline
ASCE 41-17 \cite{asce2017asce41}
& Seismic evaluation and retrofit of existing buildings.
& Provides nonlinear modelling parameters, hinge acceptance criteria, and performance levels.
& Supports beam and column hinge assignment and hinge-state evaluation in pushover analysis. \\ \hline\hline
\end{longtable}
\endgroup
\begingroup
\small
\setlength{\tabcolsep}{3pt}
\setlength{\LTpre}{6pt}
\setlength{\LTpost}{6pt}
\begin{longtable}{@{}p{3.25cm}p{3.30cm}p{3.35cm}p{2.90cm}@{}}
\caption{Summary of nonlinear modelling and seismic behaviour references reviewed}\label{tab:literature_nonlinear_modelling}\\
\hline\hline
\textbf{Author(s)} & \textbf{Study Focus} & \textbf{Methodology or Model} & \textbf{Relevance to Present Study} \\ \hline
\endfirsthead
\hline\hline
\textbf{Author(s)} & \textbf{Study Focus} & \textbf{Methodology or Model} & \textbf{Relevance to Present Study} \\ \hline
\endhead
Paulay and Priestley\newline \cite{paulay1992seismic}
& Seismic design of reinforced concrete and masonry buildings.
& Discusses ductility, plastic hinge behaviour, member deformation, and seismic design concepts.
& Supports interpretation of ductile behaviour, plastic hinge mechanism, and RC seismic response. \\ \hline
Mander et al.\newline \cite{mander1988confined}
& Stress-strain model for confined and unconfined concrete.
& Developed theoretical stress-strain relationships for confined concrete.
& Supports nonlinear material modelling of reinforced concrete members. \\ \hline
Menegotto and Pinto\newline \cite{menegotto1973cyclic}
& Cyclic behaviour of reinforced concrete plane frames.
& Proposed nonlinear steel stress-strain modelling under cyclic loading.
& Supports nonlinear modelling of reinforcement behaviour in RC members. \\ \hline
Fajfar and Gaspersic\newline \cite{fajfar1996n2}
& N2 method for seismic damage analysis of RC buildings.
& Combines pushover analysis of an MDOF system with response spectrum analysis of an equivalent SDOF system.
& Supports the theoretical basis of nonlinear static seismic assessment. \\ \hline
Fajfar\newline \cite{fajfar1999capacity}
& Capacity spectrum method based on inelastic demand spectra.
& Presents inelastic demand spectrum concept for nonlinear static procedures.
& Supports performance point and capacity spectrum interpretation. \\ \hline
Chopra and Goel\newline \cite{chopra2002modal}
& Modal pushover analysis for estimating seismic demands.
& Uses modal load patterns to improve estimation of seismic demand in buildings.
& Supports discussion of modal effects and limitations of single-mode pushover analysis. \\ \hline
Kilar and Fajfar\newline \cite{kilar1997asymmetric}
& Pushover analysis of asymmetric buildings.
& Investigates simplified pushover procedure for asymmetric structural systems.
& Supports torsional and asymmetric building response discussion. \\ \hline
Fajfar et al.\newline \cite{fajfar2005torsional}
& Torsional effects in pushover-based seismic analysis.
& Discusses torsional correction and pushover-based assessment of buildings.
& Supports the torsional response evaluation of plan-irregular buildings. \\ \hline\hline
\end{longtable}
\endgroup

\begin{figure}[!htbp]
\centering
\resizebox{0.97\textwidth}{!}{%
\begin{tikzpicture}[x=1cm,y=1cm,font=\footnotesize]
% panel 1 pushover
\begin{scope}[shift={(0,0)}]
\draw[->] (0,0)--(4.0,0) node[right]{Roof displacement};
\draw[->] (0,0)--(0,3.0) node[above]{Base shear};
\draw[thick] plot[smooth] coordinates {(0,0) (0.8,1.25) (1.6,2.05) (2.4,2.45) (3.1,2.60) (3.7,2.45)};
\node[align=center] at (2,-0.75) {Pushover capacity concept};
\end{scope}
% panel 2 hinges
\begin{scope}[shift={(5.2,0)}]
\foreach \x in {0,1.4,2.8}{\draw[thick] (\x,0)--(\x,2.6);}
\foreach \y in {0.7,1.6,2.5}{\draw[thick] (0,\y)--(2.8,\y);}
\foreach \x/\y in {0/0.7,1.4/0.7,2.8/0.7,0/1.6,1.4/1.6,2.8/1.6,0/2.5,1.4/2.5,2.8/2.5}{\fill (\x,\y) circle (0.07);}
\node[align=center] at (1.4,-0.75) {Schematic hinge locations\\at member ends};
\end{scope}
% panel 3 performance point
\begin{scope}[shift={(9.5,0)}]
\draw[->] (0,0)--(4.0,0) node[right]{Displacement};
\draw[->] (0,0)--(0,3.0) node[above]{Spectral response};
\draw[thick] plot[smooth] coordinates {(0,0) (0.8,1.4) (1.6,2.1) (2.5,2.4) (3.6,2.55)};
\draw[dashed,thick] plot[smooth] coordinates {(0.2,2.8) (1.0,2.25) (1.8,1.7) (2.7,1.15) (3.6,0.8)};
\fill (1.75,1.75) circle (0.08) node[above right]{PP};
\node[align=center] at (2,-0.75) {Capacity--demand\\performance point};
\end{scope}
\end{tikzpicture}%
}
\caption{Pushover analysis concepts.}
\label{fig:pushover_hinge_pp_concepts}
\end{figure}

The conceptual comparison in Figure~\ref{fig:pushover_hinge_pp_concepts} distinguishes three related but different nonlinear-response concepts. The capacity curve describes the global force--displacement response, hinge development represents localized member nonlinearity, and the performance point represents the demand--capacity state used for comparison; these quantities should therefore not be treated as interchangeable \cite{atc1996atc40,fema2005fema440,asce2017asce41}.

\FloatBarrier


\section{Summary of Literature Review}
The reviewed literature confirms that simplified NBC 205 provisions provide an important practical design basis for low-rise RC buildings, but their performance must be interpreted together with structural configuration, detailing quality, seismic demand, and construction practice. The Draft Final NBC 205:2012 and NBC 205:2024 documents represent distinct ready-to-use references, whereas NBC 105:2025 provides the seismic-analysis framework.

Previous studies support a range of elastic and nonlinear methods for evaluating irregular RC buildings. They also show that peak strength, performance-point response, displacement, and hinge behaviour should not be treated as interchangeable indicators. These findings support the combined interpretation of storey response, pushover capacity, performance-point response, and hinge behaviour, and the distinction between descriptive CM--CR separation and formal torsional-irregularity classification when the required diaphragm-edge displacements are unavailable.

\section{Research Gap}
Previous Nepal-focused studies provide useful evidence on MRT/RUD buildings, code revisions, and nonlinear seismic assessment, but three gaps remain relevant to this study. First, limited published evidence was identified on the response of selected plan-irregular low-rise buildings when the NBC 205:2024 ready-to-use detailing basis and the seismic-demand framework are held consistent with a Regular reference model. Second, a common seven-configuration comparison covering Regular, L, L1, L2, T, U, and Non-Parallel plans under matched soil and seismic-zone cases was not identified in the reviewed literature with storey response, CM--CR separation, pushover capacity, performance points, and hinge evidence interpreted together.

A further gap concerns geometric description. Area-based irregularity measures can describe the amount of plan area removed from a bounding rectangle, but they do not uniquely represent the location, direction, or topology of that irregularity. The reviewed literature therefore does not establish that a single area-deficit index will produce the same seismic response for different re-entrant or non-parallel plan mechanisms. In this study, the study-defined GII is compared with the computed response measures while configurations with the same GII are retained as separate cases. This allows the results to show whether equal GII values are accompanied by similar or different response, rather than assuming such a relationship in advance.

The thesis addresses these gaps through four linked stages: review of the representative Draft Final NBC 205:2012 building using the available design-check evidence; establishment of the NBC 205:2024 Regular benchmark; matched assessment of six selected plan-irregular NBC 205:2024 configurations; and synthesis of the response measures to identify practical RUD applicability boundaries. A separate four-storey Regular extension is retained as an additional height-extension case. Because a matched four-storey storey-force series equivalent to the seven-model dataset has not been verified, the extension is not used to derive a quantitative height-effect relationship. The conclusions remain limited to the analysed configurations and do not constitute a general extension of NBC 205:2024 RUD provisions.
